\documentclass[a4paper,12pt]{article}

% 页面设置
\usepackage[top=2.54cm, bottom=2.54cm, left=1.91cm, right=1.91cm]{geometry}

% 中文支持
\usepackage[UTF8]{ctex}
\usepackage{fontspec}

% 数学公式支持
\usepackage{amsmath}
\usepackage{amssymb}
\usepackage{amsfonts}
\usepackage{amsthm}

\usepackage{enumitem}

% 定理环境定义
\newtheorem*{pf}{证}
\newtheorem*{sol}{解}

% 图形支持
\usepackage{graphicx}
\usepackage{tikz}

% 超链接支持
\usepackage{hyperref}
\hypersetup{
    colorlinks=true,
    linkcolor=blue,
    filecolor=magenta,
    urlcolor=cyan,
}

% 行距设置
\linespread{1.25}

% 标题信息
\title{Mathematical Analysis II\\Week 7 Homework (April 14)}
\author{袁子强\hspace{1cm} 2025533009}
% \date{}

\begin{document}

\maketitle

\section*{Section 9.5}

10. 求下列函数在指定条件下的极值

(1) $u = x^2 + y^2$, 若 $\frac{x}{a} + \frac{y}{b} = 1$;
\begin{sol}
    构造拉格朗日函数
    $$L(x,y,\lambda)=x^2+y^2-\lambda\left(\frac{x}{a}+\frac{y}{b}-1\right)$$

    对各变量求偏导:
    $$\begin{cases}
            \dfrac{\partial L}{\partial x}=2x-\dfrac{\lambda}{a}=0 \\[8pt]
            \dfrac{\partial L}{\partial y}=2y-\dfrac{\lambda}{b}=0 \\[8pt]
            \dfrac{\partial L}{\partial \lambda}=-\dfrac{x}{a}-\dfrac{y}{b}+1=0
        \end{cases}$$

    由此可得
    $$\begin{cases}
            x=\dfrac{\lambda}{2a} \\[8pt]
            y=\dfrac{\lambda}{2b}
        \end{cases}$$

    代入约束条件,可得
    $$-\dfrac{\lambda}{2a^2}-\dfrac{\lambda}{2b^2}+1=0$$

    即
    $$\dfrac{a^2+b^2}{2a^2b^2}\lambda=1$$

    解得
    $$\lambda=\dfrac{2a^2b^2}{a^2+b^2}$$

    因此
    $$\begin{cases}
            x=\dfrac{ab^2}{a^2+b^2} \\[8pt]
            y=\dfrac{a^2b}{a^2+b^2}
        \end{cases}$$

    由于 $u$ 无极大值,因此该点为极小值点。代入可得
    $$u_{\min}=\dfrac{a^2b^2}{a^2+b^2}$$
\end{sol}

(2) $u = x + y + z$, 若 $\frac{1}{x} + \frac{1}{y} + \frac{1}{z} = 1$, $x > 0,y > 0,z > 0$;
\begin{sol}
    构造拉格朗日函数
    $$L(x,y,z,\lambda)=x+y+z-\lambda\left(\frac{1}{x}+\frac{1}{y}+\frac{1}{z}-1\right)$$

    对各变量求偏导:
    $$\begin{cases}
            \dfrac{\partial L}{\partial x}=1+\dfrac{\lambda}{x^2}=0 \\[8pt]
            \dfrac{\partial L}{\partial y}=1+\dfrac{\lambda}{y^2}=0 \\[8pt]
            \dfrac{\partial L}{\partial z}=1+\dfrac{\lambda}{z^2}=0 \\[8pt]
            \dfrac{\partial L}{\partial\lambda}=-\dfrac{1}{x}-\dfrac{1}{y}-\dfrac{1}{z}+1=0
        \end{cases}$$

    因此 $x^2=y^2=z^2=-\lambda$。由于 $x,y,z>0$,可得
    $$\dfrac{1}{x}=\dfrac{1}{y}=\dfrac{1}{z}=\dfrac{1}{3}$$

    从而 $x=y=z=3$。

    由于 $u$ 无极大值,因此该点为极小值点
    $$u_{\min}=3+3+3=9$$
\end{sol}

(3) $u = \sin x \sin y \sin z$, 若 $x + y + z = \frac{\pi}{2}$, $x > 0,y > 0,z > 0$;
\begin{sol}
    取对数得
    $$\ln u=\ln(\sin x)+\ln(\sin y)+\ln(\sin z)$$

    构造拉格朗日函数
    $$L(x,y,z,\lambda)=\ln(\sin x)+\ln(\sin y)+\ln(\sin z)-\lambda\left(x+y+z-\frac{\pi}{2}\right)$$

    对各变量求偏导:
    $$\begin{cases}
            \dfrac{\partial L}{\partial x}=\cot x-\lambda=0 \\[8pt]
            \dfrac{\partial L}{\partial y}=\cot y-\lambda=0 \\[8pt]
            \dfrac{\partial L}{\partial z}=\cot z-\lambda=0 \\[8pt]
            \dfrac{\partial L}{\partial \lambda}=-x-y-z+\dfrac{\pi}{2}=0
        \end{cases}$$

    因此 $x=y=z=\dfrac{\pi}{6}$。

    当 $x,y,z\to0^+$ 时, $u\to0$,因此该点为极大值点
    $$u_{\max}=\left(\sin\dfrac{\pi}{6}\right)^3=\dfrac{1}{8}$$
\end{sol}

(4) $u = xyz$, 若 $x + y + z = 0$ 且 $x^2 + y^2 + z^2 = 1$.
\begin{sol}
    构造拉格朗日函数
    $$L(x,y,z,\lambda,\mu)=xyz-\lambda(x+y+z)-\mu(x^2+y^2+z^2-1)$$

    对各变量求偏导:
    $$\begin{cases}
            \dfrac{\partial L}{\partial x}=yz-\lambda-2\mu x=0 \\[8pt]
            \dfrac{\partial L}{\partial y}=xz-\lambda-2\mu y=0 \\[8pt]
            \dfrac{\partial L}{\partial z}=xy-\lambda-2\mu z=0 \\[8pt]
            \dfrac{\partial L}{\partial \lambda}=-x-y-z=0      \\[8pt]
            \dfrac{\partial L}{\partial \mu}=-x^2-y^2-z^2+1=0
        \end{cases}$$

    整理可得
    $$\begin{cases}
            (z+2\mu)(x-y)=0 \\[8pt]
            (x+2\mu)(y-z)=0 \\[8pt]
            (y+2\mu)(z-x)=0
        \end{cases}$$

    若 $x=y=z$,则 $-x-y-z=0$ 与 $-x^2-y^2-z^2+1=0$ 无法同时成立,故舍去。

    因此不妨设
    $$\begin{cases}
            x=y \\
            z=-2\mu
        \end{cases}$$

    进而可得
    $$\begin{cases}
            \mu=x \\
            2x^2+4\mu^2=1
        \end{cases}$$

    解得
    $$x=\mu=\pm\dfrac{1}{\sqrt{6}}$$

    因此得到两组解
    $$\begin{cases}
            x=\dfrac{1}{\sqrt{6}} \\[8pt]
            y=\dfrac{1}{\sqrt{6}} \\[8pt]
            z=-\dfrac{2}{\sqrt{6}}
        \end{cases}\quad\text{或}\quad\begin{cases}
            x=-\dfrac{1}{\sqrt{6}} \\[8pt]
            y=-\dfrac{1}{\sqrt{6}} \\[8pt]
            z=\dfrac{2}{\sqrt{6}}
        \end{cases}$$

    从而
    $$u_{\min}=-\dfrac{1}{3\sqrt{6}},\quad u_{\max}=\dfrac{1}{3\sqrt{6}}$$
\end{sol}

11. 求下列函数在指定范围内的最大值与最小值.

(2) $z = x^2 - xy + y^2$, $\{(x,y)\mid |x| + |y| \leq 1\}$;
\begin{sol}
    首先求内部驻点:
    $$\begin{cases}
            \dfrac{\partial z}{\partial x}=2x-y=0 \\[8pt]
            \dfrac{\partial z}{\partial y}=-x+2y=0
        \end{cases}$$

    解得 $x=y=0$,此时 $z=0$。

    接下来考虑边界极值。

    \subparagraph*{情形一: $x+y=1$, $0\leq x\leq 1$}

    构造拉格朗日函数
    $$L(x,y,\lambda)=x^2-xy+y^2-\lambda(x+y-1)$$

    对各变量求偏导:
    $$\begin{cases}
            \dfrac{\partial L}{\partial x}=2x-y-\lambda=0  \\[8pt]
            \dfrac{\partial L}{\partial y}=-x+2y-\lambda=0 \\[8pt]
            \dfrac{\partial L}{\partial \lambda}=-x-y+1=0
        \end{cases}$$

    解得
    $$\begin{cases}
            x=\dfrac{1}{2} \\[8pt]
            y=\dfrac{1}{2}
        \end{cases}$$

    此时 $z=\dfrac{1}{4}$。

    \subparagraph*{情形二: $x-y=1$, $0\leq x\leq 1$}

    构造拉格朗日函数
    $$L(x,y,\lambda)=x^2-xy+y^2-\lambda(x-y-1)$$

    对各变量求偏导:
    $$\begin{cases}
            \dfrac{\partial L}{\partial x}=2x-y-\lambda=0  \\[8pt]
            \dfrac{\partial L}{\partial y}=-x+2y+\lambda=0 \\[8pt]
            \dfrac{\partial L}{\partial \lambda}=-x+y-1=0
        \end{cases}$$

    解得
    $$\begin{cases}
            x=-\dfrac{1}{2} \\[8pt]
            y=\dfrac{1}{2}
        \end{cases}$$

    该点不在定义域内,故舍去。

    \subparagraph*{情形三: $-x+y=1$, $-1\leq x\leq 0$}

    构造拉格朗日函数
    $$L(x,y,\lambda)=x^2-xy+y^2-\lambda(-x+y-1)$$

    对各变量求偏导:
    $$\begin{cases}
            \dfrac{\partial L}{\partial x}=2x-y+\lambda=0  \\[8pt]
            \dfrac{\partial L}{\partial y}=-x+2y-\lambda=0 \\[8pt]
            \dfrac{\partial L}{\partial \lambda}=x-y+1=0
        \end{cases}$$

    解得
    $$\begin{cases}
            x=\dfrac{1}{2} \\[8pt]
            y=-\dfrac{1}{2}
        \end{cases}$$

    该点不在定义域内,故舍去。

    \subparagraph*{情形四: $-x-y=1$, $-1\leq x\leq 0$}

    构造拉格朗日函数
    $$L(x,y,\lambda)=x^2-xy+y^2-\lambda(-x-y-1)$$

    对各变量求偏导:
    $$\begin{cases}
            \dfrac{\partial L}{\partial x}=2x-y+\lambda=0  \\[8pt]
            \dfrac{\partial L}{\partial y}=-x+2y+\lambda=0 \\[8pt]
            \dfrac{\partial L}{\partial \lambda}=x+y+1=0
        \end{cases}$$

    解得
    $$\begin{cases}
            x=-\dfrac{1}{2} \\[8pt]
            y=-\dfrac{1}{2}
        \end{cases}$$

    此时 $z=\dfrac{1}{4}$。

    最后计算端点处的函数值:
    $$\begin{cases}
            z(0,1)=1  \\[8pt]
            z(1,0)=1  \\[8pt]
            z(0,-1)=1 \\[8pt]
            z(-1,0)=1
        \end{cases}$$

    综上所述,
    $$z_{\min}=0,\quad z_{\max}=1$$
\end{sol}

(4) $z = x^2y(4 - x - y)$, $\{(x,y)\mid x \geq 0,y \geq 0,x + y \leq 6\}$.
\begin{sol}
    将函数展开为
    $$z=4x^2y-x^3y-x^2y^2$$

    首先求内部驻点:
    $$\begin{cases}
            \dfrac{\partial z}{\partial x}=8xy-3x^2y-2xy^2=xy(8-3x-2y)=0 \\[8pt]
            \dfrac{\partial z}{\partial y}=4x^2-2x^2y-x^3=x^2(4-2y-x)=0
        \end{cases}$$

    由于 $x>0$, $y>0$,因此
    $$\begin{cases}
            8-3x-2y=0 \\
            4-2y-x=0
        \end{cases}$$

    解得
    $$\begin{cases}
            x=2 \\
            y=1
        \end{cases}$$

    此时 $z=4$。

    接下来考虑边界极值。

    \subparagraph*{情形一: $x=0$, $0<y<6$}

    此时 $z=0$。

    \subparagraph*{情形二: $y=0$, $0<x<6$}

    此时 $z=0$。

    \subparagraph*{情形三: $x+y=6$, $0<x<6$}

    构造拉格朗日函数
    $$L(x,y,\lambda)=x^2y(4-x-y)-\lambda(x+y-6)$$

    对各变量求偏导:
    $$\begin{cases}
            \dfrac{\partial L}{\partial x}=8xy-3x^2y-2xy^2-\lambda=0 \\[8pt]
            \dfrac{\partial L}{\partial y}=4x^2-2x^2y-x^3-\lambda=0  \\[8pt]
            \dfrac{\partial L}{\partial \lambda}=-x-y+6=0
        \end{cases}$$

    解得
    $$\begin{cases}
            x=4 \\
            y=2
        \end{cases}\quad\text{或}\quad\begin{cases}
            x=0 \\
            y=6
        \end{cases}$$

    对应的函数值为 $z=-64$ 或 $z=0$。

    综上所述,
    $$z_{\min}=-64,\quad z_{\max}=4$$
\end{sol}

13. 在曲线 $\begin{cases}
        z = x^2 + 2y^2 \\
        z = 6 - 2x^2 - y^2
    \end{cases}$ 上, 求竖坐标分别为最大值和最小值的点.
\begin{sol}
    联立方程组,可得曲线在 $xy$ 平面上的投影
    $$x^2+2y^2=6-2x^2-y^2$$

    即
    $$x^2=2-y^2$$

    代入可得
    $$z=2+y^2$$

    由于 $x^2$ 有意义,因此 $0\leq y^2\leq 2$,从而
    $$2\leq z\leq 4$$

    因此
    $$z_{\min}=2,\quad \text{对应点 }(\pm\sqrt{2},0,2)$$
    $$z_{\max}=4,\quad \text{对应点 }(0,\pm\sqrt{2},4)$$
\end{sol}

16. 已知平行六面体所有各棱长之和为 $12a$, 求其最大体积.
\begin{sol}
    固定棱长时,当平行六面体为长方体时体积最大。因此问题转化为求各棱长之和为 $12a$ 的长方体的最大体积。

    设长、宽、高分别为 $x,y,z$,则体积为
    $$V=xyz$$

    约束条件为
    $$4x+4y+4z=12a$$

    即
    $$x+y+z=3a$$

    构造拉格朗日函数
    $$L(x,y,z,\lambda)=xyz-\lambda(x+y+z-3a)$$

    对各变量求偏导:
    $$\begin{cases}
            \dfrac{\partial L}{\partial x}=yz-\lambda=0 \\[8pt]
            \dfrac{\partial L}{\partial y}=xz-\lambda=0 \\[8pt]
            \dfrac{\partial L}{\partial z}=xy-\lambda=0 \\[8pt]
            \dfrac{\partial L}{\partial \lambda}=-x-y-z+3a=0
        \end{cases}$$

    解得 $x=y=z=a$,因此最大体积为
    $$V_{\max}=a^3$$
\end{sol}

19. 椭球体 $\frac{x^2}{a^2} + \frac{y^2}{b^2} + \frac{z^2}{c^2} \leq 1$ 的内接长方体中, 求体积最大的长方体的体积.
\begin{sol}
    设在第一卦限内的顶点坐标为 $(x,y,z)$,则体积为
    $$V=8xyz$$

    约束条件为
    $$\dfrac{x^2}{a^2} + \dfrac{y^2}{b^2} + \dfrac{z^2}{c^2} = 1$$

    构造拉格朗日函数
    $$L(x,y,z,\lambda)=8xyz-\lambda\left(\dfrac{x^2}{a^2} + \dfrac{y^2}{b^2} + \dfrac{z^2}{c^2} - 1\right)$$

    对各变量求偏导:
    $$\begin{cases}
            \dfrac{\partial L}{\partial x}=8yz-\lambda\dfrac{2x}{a^2}=0 \\[8pt]
            \dfrac{\partial L}{\partial y}=8xz-\lambda\dfrac{2y}{b^2}=0 \\[8pt]
            \dfrac{\partial L}{\partial z}=8xy-\lambda\dfrac{2z}{c^2}=0 \\[8pt]
            \dfrac{\partial L}{\partial \lambda}=-\dfrac{x^2}{a^2} - \dfrac{y^2}{b^2} - \dfrac{z^2}{c^2} + 1=0
        \end{cases}$$

    由此可得
    $$\dfrac{x^2}{a^2} = \dfrac{y^2}{b^2} = \dfrac{z^2}{c^2} = \dfrac{1}{3}$$

    因此
    $$x=\dfrac{a\sqrt{3}}{3},\quad y=\dfrac{b\sqrt{3}}{3},\quad z=\dfrac{c\sqrt{3}}{3}$$

    从而最大体积为
    $$V_{\max}=\dfrac{8\sqrt{3}}{9}abc$$
\end{sol}

21. 设曲面 $S$: $\sqrt{x} + \sqrt{y} + \sqrt{z} = \sqrt{a}\ (a > 0)$

(1) 证明 $S$ 上任意点处的切平面与各坐标轴的截距之和等于 $a$.
\begin{pf}
    设函数
    $$F(x,y,z)=\sqrt{x} + \sqrt{y} + \sqrt{z} - \sqrt{a}=0$$

    对各变量求偏导:
    $$\begin{cases}
            \dfrac{\partial F}{\partial x}=\dfrac{1}{2\sqrt{x}} \\[8pt]
            \dfrac{\partial F}{\partial y}=\dfrac{1}{2\sqrt{y}} \\[8pt]
            \dfrac{\partial F}{\partial z}=\dfrac{1}{2\sqrt{z}}
        \end{cases}$$

    因此在切点 $(x_0,y_0,z_0)$ (满足 $\sqrt{x_0} + \sqrt{y_0} + \sqrt{z_0} = \sqrt{a}$) 处的切平面法向量为
    $$\left(\dfrac{1}{\sqrt{x_0}}, \dfrac{1}{\sqrt{y_0}}, \dfrac{1}{\sqrt{z_0}}\right)$$

    切平面方程为
    $$\dfrac{x-x_0}{\sqrt{x_0}}+\dfrac{y-y_0}{\sqrt{y_0}}+\dfrac{z-z_0}{\sqrt{z_0}}=0$$

    即
    $$\dfrac{x}{\sqrt{x_0}}+\dfrac{y}{\sqrt{y_0}}+\dfrac{z}{\sqrt{z_0}}=\sqrt{a}$$

    因此截距分别为 $\sqrt{ax_0}$, $\sqrt{ay_0}$, $\sqrt{az_0}$,求和可得
    $$\sqrt{a}(\sqrt{x_0}+\sqrt{y_0}+\sqrt{z_0})=\sqrt{a} \cdot \sqrt{a}=a$$

    Q.E.D.
\end{pf}

(2) 在 $S$ 上求一切平面, 使此切平面与三坐标面所围成的四面体体积最大, 并求四面体体积的最大值.
\begin{sol}
    设切点坐标为 $(x,y,z)$。由(1)可知,切平面在各坐标轴上的截距分别为 $\sqrt{ax}$, $\sqrt{ay}$, $\sqrt{az}$,因此四面体体积为
    $$V=\dfrac{1}{6}\sqrt{ax}\cdot\sqrt{ay}\cdot\sqrt{az}=\dfrac{a^{\frac{3}{2}}\sqrt{xyz}}{6}$$

    约束条件为
    $$\sqrt{x} + \sqrt{y} + \sqrt{z} = \sqrt{a}$$

    令 $u=\sqrt{x}$, $v=\sqrt{y}$, $w=\sqrt{z}$,则问题转化为在约束条件
    $$u+v+w=\sqrt{a}$$

    下求
    $$V=\dfrac{a^{\frac{3}{2}}uvw}{6}$$

    的最大值。显然当 $u=v=w=\dfrac{\sqrt{a}}{3}$ 时体积最大,此时
    $$x=y=z=\dfrac{a}{9}$$

    因此最大体积为
    $$V_{\max}=\dfrac{a^3}{162}$$
\end{sol}

\section*{Section 9.6}

2. 设 $\boldsymbol{\omega} = \omega_1 \boldsymbol{i} + \omega_2 \boldsymbol{j} + \omega_3 \boldsymbol{k}$ 是一个常值向量, 求向量场
$$
    \boldsymbol{v} = \boldsymbol{\omega} \times \boldsymbol{r}
$$
的旋度 $\nabla \times \boldsymbol{v}$, 并给出合理的物理解释. 这里 $\boldsymbol{r} = x\boldsymbol{i} + y\boldsymbol{j} + z\boldsymbol{k}$ 是位置向量.
\begin{sol}
    首先计算向量场 $\boldsymbol{v}$:
    $$\boldsymbol{v}=\boldsymbol{\omega}\times \boldsymbol{r}=\begin{vmatrix}
            \boldsymbol{i} & \boldsymbol{j} & \boldsymbol{k} \\
            \omega_1       & \omega_2       & \omega_3       \\
            x              & y              & z
        \end{vmatrix}=(\omega_2z-\omega_3y,\omega_3x-\omega_1z,\omega_1y-\omega_2x)$$

    进而计算旋度:
    $$\nabla\times \boldsymbol{v}=\begin{vmatrix}
            \boldsymbol{i}      & \boldsymbol{j}      & \boldsymbol{k}      \\
            \dfrac{\partial}{\partial x}          & \dfrac{\partial}{\partial y}          & \dfrac{\partial}{\partial z}          \\
            \omega_2z-\omega_3y & \omega_3x-\omega_1z & \omega_1y-\omega_2x
        \end{vmatrix}=(2\omega_1,2\omega_2,2\omega_3)=2\boldsymbol{\omega}$$

    即
    $$\nabla\times \boldsymbol{v}=\nabla\times(\boldsymbol{\omega}\times \boldsymbol{r})=2\boldsymbol{\omega}$$

    物理解释: $\boldsymbol{v}$ 表示刚体绕定轴以角速度 $\boldsymbol{\omega}$ 旋转时的线速度场,其旋度等于角速度的 $2$ 倍。
\end{sol}


4. 设 $\boldsymbol{w}$ 是常向量, $\boldsymbol{r} = x\boldsymbol{i} + y\boldsymbol{j} + z\boldsymbol{k}$, $r = |\boldsymbol{r}|$, 求

(1) $\operatorname{div}\left[(\boldsymbol{r} \cdot \boldsymbol{w})\boldsymbol{w}\right]$;
\begin{sol}
    由乘积法则可得
    $$\nabla\cdot[(\boldsymbol{r}\cdot \boldsymbol{w})\boldsymbol{w}]=\nabla(\boldsymbol{r}\cdot \boldsymbol{w})\cdot \boldsymbol{w}+(\boldsymbol{r}\cdot \boldsymbol{w})\nabla\cdot \boldsymbol{w}$$

    由于 $\boldsymbol{w}$ 是常向量,因此 $\nabla\cdot \boldsymbol{w}=0$,从而
    $$\nabla\cdot[(\boldsymbol{r}\cdot \boldsymbol{w})\boldsymbol{w}]=\nabla(\boldsymbol{r}\cdot \boldsymbol{w})\cdot \boldsymbol{w}=\boldsymbol{w}\cdot \boldsymbol{w}=|\boldsymbol{w}|^2$$
\end{sol}

(2) $\operatorname{div} \frac{\boldsymbol{r}}{r}$;
\begin{sol}
    由散度公式可得
    $$\nabla\cdot\dfrac{\boldsymbol{r}}{r}=\nabla\cdot\left(\boldsymbol{r}\cdot\dfrac{1}{r}\right)=\dfrac{\nabla\cdot\boldsymbol{r}}{r}+\boldsymbol{r}\cdot\nabla\dfrac{1}{r}$$

    其中 $\nabla\cdot\boldsymbol{r}=3$,且
    $$\nabla\dfrac{1}{r}=-\dfrac{1}{r^2}\nabla r=-\dfrac{1}{r^2}\cdot\dfrac{\boldsymbol{r}}{r}=-\dfrac{\boldsymbol{r}}{r^3}$$

    因此
    $$\nabla\cdot\dfrac{\boldsymbol{r}}{r}=\dfrac{3}{r}+\boldsymbol{r}\cdot\left(-\dfrac{\boldsymbol{r}}{r^3}\right)=\dfrac{3}{r}-\dfrac{r^2}{r^3}=\dfrac{2}{r}$$
\end{sol}

(3) $\operatorname{div} \left(\boldsymbol{w} \times \boldsymbol{r}\right)$;
\begin{sol}
    由向量恒等式可得
    $$\nabla\cdot(\boldsymbol{w}\times \boldsymbol{r})=\boldsymbol{r}\cdot(\nabla\times \boldsymbol{w})-\boldsymbol{w}\cdot(\nabla\times \boldsymbol{r})$$

    由于 $\boldsymbol{w}$ 是常向量,因此 $\nabla\times \boldsymbol{w}=\boldsymbol{0}$。又因为
    $$\nabla\times \boldsymbol{r}=\begin{vmatrix}
            \boldsymbol{i} & \boldsymbol{j} & \boldsymbol{k} \\
            \dfrac{\partial}{\partial x} & \dfrac{\partial}{\partial y} & \dfrac{\partial}{\partial z} \\
            x & y & z
        \end{vmatrix}=\boldsymbol{0}$$

    因此
    $$\nabla\cdot(\boldsymbol{w}\times \boldsymbol{r})=0$$
\end{sol}

(4) $\operatorname{div} \left(r^2 \boldsymbol{w}\right)$.
\begin{sol}
    由乘积法则可得
    $$\nabla\cdot(r^2\boldsymbol{w})=\nabla r^2\cdot \boldsymbol{w}+r^2\nabla\cdot \boldsymbol{w}$$

    由于 $\boldsymbol{w}$ 是常向量,因此 $\nabla\cdot \boldsymbol{w}=0$,从而
    $$\nabla\cdot(r^2\boldsymbol{w})=\nabla r^2\cdot \boldsymbol{w}$$

    其中 $\nabla r^2=2r\nabla r=2r\cdot\dfrac{\boldsymbol{r}}{r}=2\boldsymbol{r}$,因此
    $$\nabla\cdot(r^2\boldsymbol{w})=2\boldsymbol{r}\cdot \boldsymbol{w}$$
\end{sol}


6. 设 $\boldsymbol{w}$ 是常向量, $\boldsymbol{r} = x\boldsymbol{i} + y\boldsymbol{j} + z\boldsymbol{k}$, $r = |\boldsymbol{r}|$, 求:

(1) $\operatorname{rot}\left(\boldsymbol{w} \times \boldsymbol{r}\right)$;
\begin{sol}
    由向量恒等式可得
    $$\nabla\times(\boldsymbol{w}\times\boldsymbol{r})=(\boldsymbol{r}\cdot\nabla)\boldsymbol{w}-(\boldsymbol{w}\cdot\nabla)\boldsymbol{r}+\boldsymbol{w}(\nabla\cdot\boldsymbol{r})-\boldsymbol{r}(\nabla\cdot\boldsymbol{w})$$

    由于 $\boldsymbol{w}$ 是常向量,因此 $(\boldsymbol{r}\cdot\nabla)\boldsymbol{w}=\boldsymbol{0}$, $\nabla\cdot\boldsymbol{w}=0$。又因为
    $$(\boldsymbol{w}\cdot\nabla)\boldsymbol{r}=\boldsymbol{w},\quad \nabla\cdot\boldsymbol{r}=3$$

    因此
    $$\nabla\times(\boldsymbol{w}\times\boldsymbol{r})=\boldsymbol{0}-\boldsymbol{w}+3\boldsymbol{w}-\boldsymbol{0}=2\boldsymbol{w}$$
\end{sol}

(2) $\operatorname{rot}\left[f(r)\boldsymbol{r}\right]$;
\begin{sol}
    由旋度公式可得
    $$\nabla\times[f(r)\boldsymbol{r}]
        =\nabla f(r)\times\boldsymbol{r}+f(r)(\nabla\times\boldsymbol{r})$$

    其中
    $$\nabla f(r)=f'(r)\nabla r=f'(r)\dfrac{\boldsymbol{r}}{r},\quad \nabla\times\boldsymbol{r}=\boldsymbol{0}$$

    因此
    $$\nabla\times[f(r)\boldsymbol{r}]=f'(r)\dfrac{\boldsymbol{r}}{r}\times\boldsymbol{r}=\boldsymbol{0}$$
\end{sol}

(3) $\operatorname{rot}\left[f(r)\boldsymbol{w}\right]$;
\begin{sol}
    由旋度公式可得
    $$\nabla\times[f(r)\boldsymbol{w}]
        =\nabla f(r)\times\boldsymbol{w}+f(r)(\nabla\times\boldsymbol{w})$$

    由于 $\boldsymbol{w}$ 是常向量,因此 $\nabla\times\boldsymbol{w}=\boldsymbol{0}$,从而
    $$\nabla\times[f(r)\boldsymbol{w}]=\nabla f(r)\times\boldsymbol{w}=f'(r)\dfrac{\boldsymbol{r}}{r}\times\boldsymbol{w}$$
\end{sol}

(4) $\operatorname{div}\left[\boldsymbol{r} \times f(r)\boldsymbol{w}\right]$.
\begin{sol}
    由散度公式可得
    $$\nabla\cdot[\boldsymbol{r} \times f(r)\boldsymbol{w}]=f(r)\boldsymbol{w}\cdot(\nabla\times\boldsymbol{r})-\boldsymbol{r}\cdot(\nabla\times f(r)\boldsymbol{w})$$

    由于 $\nabla\times\boldsymbol{r}=\boldsymbol{0}$,且由第(3)小题可知 $\nabla\times f(r)\boldsymbol{w}=f'(r)\dfrac{\boldsymbol{r}}{r}\times\boldsymbol{w}$,因此
    $$\nabla\cdot[\boldsymbol{r} \times f(r)\boldsymbol{w}]=-\boldsymbol{r}\cdot\left(f'(r)\dfrac{\boldsymbol{r}}{r}\times\boldsymbol{w}\right)$$

    由于 $\boldsymbol{r}$ 与 $\boldsymbol{r}\times\boldsymbol{w}$ 垂直,因此 $\boldsymbol{r}\cdot\left(f'(r)\dfrac{\boldsymbol{r}}{r}\times\boldsymbol{w}\right)=0$,从而
    $$\nabla\cdot[\boldsymbol{r} \times f(r)\boldsymbol{w}]=0$$
\end{sol}


7. 设 $\boldsymbol{w}$ 是常向量, $\boldsymbol{r} = x\boldsymbol{i} + y\boldsymbol{j} + z\boldsymbol{k}$, $r = |\boldsymbol{r}|$, $f(r)$ 是 $r$ 的可微函数, 试通过 $\nabla$ 运算求:

(1) $\nabla\left(\boldsymbol{w} \cdot f(r)\boldsymbol{r}\right)$;
\begin{sol}
    由梯度公式可得
    $$\nabla(\boldsymbol{w}\cdot f(r)\boldsymbol{r})=\boldsymbol{w}\times(\nabla\times(f(r)\boldsymbol{r}))+(f(r)\boldsymbol{r})\times(\nabla\times\boldsymbol{w})+(\boldsymbol{w}\cdot\nabla)(f(r)\boldsymbol{r})+(f(r)\boldsymbol{r}\cdot\nabla)\boldsymbol{w}$$

    由于 $\boldsymbol{w}$ 是常向量,因此 $\nabla\times\boldsymbol{w}=\boldsymbol{0}$, $(f(r)\boldsymbol{r}\cdot\nabla)\boldsymbol{w}=\boldsymbol{0}$。由第6题(2)可知 $\nabla\times(f(r)\boldsymbol{r})=\boldsymbol{0}$,因此
    $$\nabla(\boldsymbol{w}\cdot f(r)\boldsymbol{r})=(\boldsymbol{w}\cdot\nabla)(f(r)\boldsymbol{r})$$

    展开可得
    $$(\boldsymbol{w}\cdot\nabla)(f(r)\boldsymbol{r})=f(r)(\boldsymbol{w}\cdot\nabla)\boldsymbol{r}+(\boldsymbol{w}\cdot\nabla f(r))\boldsymbol{r}$$

    其中 $(\boldsymbol{w}\cdot\nabla)\boldsymbol{r}=\boldsymbol{w}$,且
    $$\boldsymbol{w}\cdot\nabla f(r)=f'(r)\boldsymbol{w}\cdot\dfrac{\boldsymbol{r}}{r}$$

    因此
    $$\nabla(\boldsymbol{w}\cdot f(r)\boldsymbol{r})=f(r)\boldsymbol{w}+f'(r)\dfrac{\boldsymbol{w}\cdot\boldsymbol{r}}{r}\boldsymbol{r}$$
\end{sol}

(2) $\nabla \cdot \left(\boldsymbol{w} \times f(r)\boldsymbol{r}\right)$;
\begin{sol}
    由散度公式可得
    $$\nabla\cdot[\boldsymbol{w} \times f(r)\boldsymbol{r}]=f(r)\boldsymbol{r}\cdot(\nabla\times\boldsymbol{w})-\boldsymbol{w}\cdot(\nabla\times f(r)\boldsymbol{r})$$

    由于 $\nabla\times\boldsymbol{w}=\boldsymbol{0}$,且由第6题(2)可知 $\nabla\times f(r)\boldsymbol{r}=\boldsymbol{0}$,因此
    $$\nabla\cdot[\boldsymbol{w} \times f(r)\boldsymbol{r}]=0$$
\end{sol}

(3) $\nabla \times \left(\boldsymbol{w} \times f(r)\boldsymbol{r}\right)$.
\begin{sol}
    由旋度公式可得
    $$\nabla\times(\boldsymbol{w}\times f(r)\boldsymbol{r})=(f(r)\boldsymbol{r}\cdot\nabla)\boldsymbol{w}-(\boldsymbol{w}\cdot\nabla)(f(r)\boldsymbol{r})+\boldsymbol{w}(\nabla\cdot(f(r)\boldsymbol{r}))-f(r)\boldsymbol{r}(\nabla\cdot\boldsymbol{w})$$

    由于 $\boldsymbol{w}$ 是常向量,因此 $(f(r)\boldsymbol{r}\cdot\nabla)\boldsymbol{w}=\boldsymbol{0}$, $\nabla\cdot\boldsymbol{w}=0$,从而
    $$\nabla\times(\boldsymbol{w}\times f(r)\boldsymbol{r})=-(\boldsymbol{w}\cdot\nabla)(f(r)\boldsymbol{r})+\boldsymbol{w}(\nabla\cdot(f(r)\boldsymbol{r}))$$

    其中由第(1)小题可知
    $$(\boldsymbol{w}\cdot\nabla)(f(r)\boldsymbol{r})=f(r)\boldsymbol{w}+f'(r)\dfrac{\boldsymbol{w}\cdot\boldsymbol{r}}{r}\boldsymbol{r}$$

    且
    $$\nabla\cdot(f(r)\boldsymbol{r})=f(r)\nabla\cdot\boldsymbol{r}+\boldsymbol{r}\cdot\nabla f(r)=3f(r)+\boldsymbol{r}\cdot f'(r)\dfrac{\boldsymbol{r}}{r}=3f(r)+rf'(r)$$

    因此
    $$\begin{aligned}
        \nabla\times(\boldsymbol{w}\times f(r)\boldsymbol{r})
        &=-(f(r)\boldsymbol{w}+f'(r)\dfrac{\boldsymbol{w}\cdot\boldsymbol{r}}{r}\boldsymbol{r})+\boldsymbol{w}(3f(r)+rf'(r)) \\
        &=(3f(r)+rf'(r))\boldsymbol{w}-f(r)\boldsymbol{w}-f'(r)\dfrac{\boldsymbol{w}\cdot\boldsymbol{r}}{r}\boldsymbol{r} \\
        &=(2f(r)+rf'(r))\boldsymbol{w}-f'(r)\dfrac{\boldsymbol{w}\cdot\boldsymbol{r}}{r}\boldsymbol{r}
    \end{aligned}$$
\end{sol}

\end{document}